\section{地形—载荷—能量统一计算模型}\label{sec:physics}
四问要在同一标准场景下衔接，首先必须让同一条路线得到相同的时间、能耗和状态序列。本文将运输架次表示为“货箱集合、机型、访问顺序”，公共引擎据此返回逐段载荷、飞行阶段、交付偏移、总时长与总能耗。上层算法决定怎样组批和排程，不能改变这套物理计算。

\subsection{DEM 航高与飞行阶段}
山区航线的代价并不只由水平距离决定：沿途一处高地就会抬升整个航段的计划巡航海拔。记线路经过的闭 DEM 像元集合为 $\mathcal C_{ij}$，像元高程为 $h_c$；O01 作业海拔取地面高程，服务区作业海拔取地面以上 30 m。对不同节点间航段，有
\begin{align}
H_{ij}&=\max_{c\in\mathcal C_{ij}}h_c+50,\qquad
h^+_{ij}=H_{ij}-z_i,\quad h^-_{ij}=H_{ij}-z_j,\label{eq:altitude}\\
\tau_{gij}&=\frac{h^+_{ij}}{v_g^\uparrow}+\frac{d_{ij}}{v_g^c}+\frac{h^-_{ij}}{v_g^\downarrow}.\label{eq:flight}
\end{align}
长度以 m、速度以 m/s 计，所得时间为 s。每次完成服务区交付后，下一航段从该点作业高度重新爬升；不能把多点配送当成一次起降、全程等高巡航。同一节点到自身的零航段则取零距离、零升降。
\paperfig{fig03_physics}{同一水平航段的运输代价仍随地形航高和载荷改变，空载返程也消耗能源。S014 的真实 DEM 剖面对应上图，实际 Q1 批次的去程载货与返程空载能量分解对应下图；下降附加能耗为零遵循题设规则。}{physics}
图\ref{fig:physics}把航高与能量联系起来：终点较低并不意味着沿途无需高爬升，空载返程也不能省略。因而相同水平距离可对应不同的山区运输代价，后续载荷边界须逐航段计算。

\subsection{动态载荷与等效航程}
载荷相关能耗会改变配送路线的可行域\cite{dorling}，但不同能耗模型具有各自适用条件\cite{zhang}。本题已有统一计算规则，故直接使用官方等效航程，不再用外部经验模型替换。机型 $g$ 携带有效载荷 $q$ 时，
\begin{equation}
L_g(q)=L_g^0-(L_g^0-L_g^F)(q/Q_g)^{3/2},\qquad 0\le q\le Q_g.\label{eq:range}
\end{equation}
对架次 $p$，令 $\mathcal B_p$ 为其货箱集合，$\pi_p$ 为从 O01 出发并返回的节点序列。进入航段 $i\to j$ 时，
\begin{equation}
q_{p,ij}=\sum_{b\in\mathcal B_p:\,b\text{尚未完成交付}}m_b.\label{eq:payloadstate}
\end{equation}
每个服务区完成交接后移除对应货箱，再计算下一段。由此，载荷不是“起飞质量一直保持不变”，也不是“每段都按空载”；全部货物交付后的返航段才取 $q=0$。

\subsection{分阶段能量与返航储备}
水平航程和正向爬升分别消耗电量：
\begin{equation}
E_{gij}(q)=E_g\frac{d_{ij}}{L_g(q)}+
\frac{(m_g^0+q)g_0h^+_{ij}}{\eta_g\,3.6\times10^6}.\label{eq:energy}
\end{equation}
其中 $m_g^0$ 为含电池空载总质量，$g_0=9.80665$ m/s$^2$，$\eta_g$ 为爬升效率。第二项由焦耳换算为 kWh，与第一项一致。下降附加项按题设取零，并不意味着真实下降过程不耗电，而是本题等效模型的规定。

架次可行性同时要求
\begin{equation}
\begin{gathered}
\sum_{b\in\mathcal B_p}m_b\le Q_{g_p},\qquad
\sum_{b\in\mathcal B_p}v_b\le V_{g_p},\\
E_p=\sum_{(i,j)\in\pi_p}E_{g_pij}(q_{p,ij})
\le(1-\rho_{g_p})E_{g_p}.\label{eq:feasible}
\end{gathered}
\end{equation}
质量限制回答“是否超载”，体积限制回答“能否装入”，能源限制回答“能否完成全程并留出返航储备”。由于只卸货、不取货，起飞质量和体积是全程最大值；电量则必须累计所有航段。共享这一判据，比单纯要求四问使用相同参数名称更重要。

\subsection{无人机与能源资源的独立占用}
作业从固定准备和逐箱装载开始。各段飞行之后，按服务区基础交接时间及逐箱附加时间计算完成时刻；货箱的交付偏移 $\delta_{pb}$ 取交接结束，不能用到达代替。公共引擎最终给出 $(t_p,E_p,\delta_{pb})$，Q2、Q3 再通过 $C_b=s_p+\delta_{pb}$ 转成绝对时刻。

无人机返航即可结束该运输架次占用；电池还须充至满电。令 $\sigma_p=1-E_p/E_{g_p}$，官方两阶段充电时间为
\begin{equation}
 t_{\rm chg}(\sigma)=T_{\rm full}
 \begin{cases}
 0.65(0.90-\sigma)/0.90+0.35,&0\le\sigma<0.90,\\
 0.35(1-\sigma)/0.10,&0.90\le\sigma\le1.
 \end{cases}\label{eq:charge}
\end{equation}
于是无人机和电池占用区间分别为 $[s_p,s_p+t_p)$ 与 $[s_p,s_p+t_p+t_{\rm chg}(\sigma_p))$。半开区间允许资源在前一任务释放的时刻立即投入下一任务。它们将在 Q2 约束排程，在 Q3 约束时序 LP，并在 Q4 转为独立资源需求，构成四问之间可核验的公共接口。
\FloatBarrier
